\documentclass[11pt]{article}
\usepackage{fontspec}            % (requiere XeLaTeX/LuaLaTeX)
\setmainfont{Fira Sans}          % Usa Fira Sans como fuente del texto
\setsansfont{Fira Sans}          % (opcional) por si usas \sffamily
\setmonofont{Fira Mono}          % (opcional) para código/monoespaciada
\usepackage{setspace}    
\onehalfspacing
\usepackage{graphicx}    
\usepackage{amsmath}     
%\usepackage{natbib}      % Apagado en esta tarea
\usepackage{array}       
\usepackage{multirow}    
\usepackage{siunitx}     
\usepackage{textcomp}    
%\usepackage[utf8]{inputenc}  %ignorado, el compilador ya tiene utf8
\usepackage[spanish,es-tabla]{babel}  
\usepackage{csquotes}        
\usepackage[top=1in,right=1in,left=1in,bottom=1in]{geometry} 
\usepackage{makecell}
\usepackage{hyperref}        

\hypersetup{
    unicode=true,            
    pdftoolbar=true,         
    pdfmenubar=true,         
    pdffitwindow=false,      
    pdfstartview={FitH},     
    pdftitle={},             
    pdfauthor={},            
    pdfsubject={},           
    pdfcreator={},           
    pdfproducer={},          
    pdfkeywords={},          
    pdfnewwindow=true,       
    colorlinks=true,         
    linkcolor=blue,          
    citecolor=blue,          
    filecolor=blue,          
    urlcolor=blue            
}
%% DECLARACIONES BIBLIOGRÁFICAS
\usepackage[backend=biber,style=apa,natbib=true,language=spanish]{biblatex}
\DeclareBibliographyDriver{newspaper}{
  \usebibmacro{author}%
  \newunit\newblock
  % Fecha en español: (dd de mm de aaaa)
  \printtext[parens]{%
    \iffieldsequal{year}{}
      {}
      {%
        \printfield[day]{day}\space
        de\space
        \printfield[month]{month}\space
        de\space
        \printfield{year}%
      }%
  }%
  \newunit\newblock
  % Título
  \printfield{title}\addperiod\newunit\newblock
  % Periódico en cursiva
  \printfield[journaltitle]{journaltitle}\addperiod\newunit\newblock
  % URL
  \printfield{url}%
  \finentry
}
\addbibresource{Referencias.bib}

\setlength\bibitemsep{1.0\baselineskip}
\DeclareLanguageMapping{spanish}{spanish-apa}
% --- APA en español: "y" en vez de "&" ---
\DefineBibliographyStrings{spanish}{
  and       = {y},
  andothers = {et\space al\adddot} % o {y\space otros} si lo prefieres
}

% Citas en el texto: \citep, \parencite, \textcite
\DeclareDelimFormat[parencite]{finalnamedelim}{\addspace\bibstring{and}\space}
\DeclareDelimFormat[textcite]{finalnamedelim}{\addspace\bibstring{and}\space}

% Bibliografía (APA usa un delimitador propio en orden apellido-nombre)
\DeclareDelimFormat[bib]{finalnamedelim:apa:family-given}{%
  \finalandcomma\addspace\bibstring{and}\space}

% Personalización para agregar bullets a la bibliografía
\defbibenvironment{bibliography}
  {\list{\labelitemi} % Usa el símbolo de bullet por defecto
    {\setlength{\leftmargin}{0pt}%
     \setlength{\itemindent}{\labelwidth}%
     \addtolength{\itemindent}{\labelsep}%
     \setlength{\itemsep}{\bibitemsep}%
     \setlength{\parsep}{\bibparsep}}}
  {\endlist}
  {\item}
%% FIN DE LA PERSONALIZACIÓN DE BIBLIOGRAFÍA A APA ESPAÑOL ---

\title{SOCI 4186-002\\ Tarea de Proyecto de Investigación \textnumero 4 \\ Propuesta de investigación \\ \textsc{Individual o Nombre de Grupo}}
\author{Pon tu nombre aquí \\ Número de estudiante}
\date{\today} %dejen esto como está, con \today

\begin{document}
\singlespacing
\maketitle
Fecha de entrega: 31 de octubre de 2025 a las 23:59 vía Teams.

\begin{enumerate}
    \item Descargarán este documento en el GitHub de la clase y lo subirán a la una carpeta en Overleaf. El compilador a usar es \texttt{XeLaTeX}.
    \item Modificarán la N de SOCI 4186-00N según su sección. Si es individual, dejen \textsc{Individual}; si es grupal, cambien el texto de plantilla por el nombre del grupo.
    \item Cambien el campo \texttt{\textbackslash author\{Pon tu nombre aquí\}}. Si es individual: su nombre completo seguido de \textbackslash\textbackslash y su número de estudiante. Si es en grupo, usen el formato: \texttt{\textbackslash author\{Estudiante 1 \textbackslash\textbackslash \ \textnumero\ de Estudiante 1 \textbackslash and Estudiante 2 \textbackslash\textbackslash \ \textnumero\ de Estudiante 2\textbackslash and Estudiante 3 \textbackslash\textbackslash \ \textnumero\ de Estudiante 3 \textbackslash and Estudiante 4 \textbackslash\textbackslash\ \textnumero\ de Estudiante 4\}}.
    \item Presten atención a la ortografía y sintaxis española (acentos, tildes, diéresis) o escriban en una de las tres modalidades de inglés discutidas: británico, Oxford o estadounidense. Cíñanse a una lengua a través de todo el documento.
%    \item Completarán como mejor sea posible la tarea, y podrán, si entienden necesario hacerlo, trabajar la tarea con compañeres de clase. Si así lo hicieren, pondrán antes de la sección 1 una oración indicándome con quién(es) trabajó la tarea. De lo contrario, dejarán la línea que está actualmente escrita.
    \item Descarguen el PDF tras compilar (icono de flecha descendente) y lo someterán a Teams.
\end{enumerate}

\section*{Formato del documento}
Siga esta plantilla, borrando las instrucciones si así deseare tras concluir las respuestas o tener el PDF de instrucciones a la mano. 
\onehalfspacing
\section*{Propuesta de investigación}
Elaboración de la propuesta de investigación. El documento podrá ser entregado incluyendo en sí las siguientes partes:

\begin{enumerate}
    \item[1.º] Su pregunta de investigación, modificada si ha sido necesario alterarla.
    \item[2.º] Identificar el método o métodos que se propone utilizar para estudiar su tema.
    \item[3.º] Especificar cuál será su unidad de análisis en su investigación.
    \item[4.º] Identificar cuál es la población objetivo (¿a quiénes o qué entidades quieren generalizar o entender?) y qué método de muestreo (probabilístico o no probabilístico, cuál) usará.
    \item[5.º] Desarrollar el plan para recolectar sus datos (vean el apéndice, que enviara hace unas semanas y reenvío junto a este documento para las instrucciones específicas de cada método).
    \item[6.º] Provea un plan que indique cómo se asegurará que su investigación se conformará a estándares de ética:
    \begin{enumerate}
        \item Explique cómo su investigación se adhiere a los principios Belmont: respeto por las personas, beneficencia y justicia.
        \item Determine y reporte si su método de investigación requiere aprobación de CIPSHI (el profesor se encargará de someter un formulario global).
    \end{enumerate}
\end{enumerate}



Si entiende necesario, puede incluir citas sobre este tema usando la bibliografía que ha ido construyendo en trabajos previos. En ese caso, use adecuadamente citas textuales --\citet{hall2001introduction}-- o parentéticas --\citep[e.g.,][, entre otros]{hall2009institutional}.

\printbibliography[title={Bibliografía}, heading=subbibliography]

\end{document}